% ----------------------------------
%   Black-Scholes (Chapter 4)
% ----------------------------------

\begin{frame}
\frametitle{Black-Scholes. Dynamic replication}
\begin{scriptsize}
\begin{itemize}
\item Arbitrage-free market, constant rate $r$, continuous rebalancing.
The underlying and the bond follow
\begin{equation*}
    dS_t = \mu S_t\, dt + \sigma S_t\, dW_t,
    \qquad
    d\beta_t = r \beta_t\, dt.
\end{equation*}
\item Replicating portfolio $V_t = a_t S_t + b_t \beta_t$ with $V_T = (S_T - K)^+$,
and \textbf{self-financing}
\begin{equation*}
    dV_t = a_t\, dS_t + b_t\, d\beta_t .
\end{equation*}
\item Writing $V_t = f(t, S_t)$, applying It\^{o}'s formula and matching the $dW_t$ and $dt$ terms,
\begin{equation*}
    a_t = \frac{\partial f}{\partial S} \quad \text{(delta hedge)},
\end{equation*}
which gives the \textbf{Black-Scholes PDE}
\begin{equation*}
    \frac{\partial f}{\partial t} =
        r f
        - r S \frac{\partial f}{\partial S}
        - \frac{1}{2} \sigma^2 S^2 \frac{\partial^2 f}{\partial S^2},
    \qquad f(T, S) = (S - K)^+ .
\end{equation*}
\item The drift $\mu$ does not appear: the price does not depend on the investors' view of the asset.
\end{itemize}
\smallskip
Black, F., Scholes, M. (1973). \emph{The Pricing of Options and Corporate Liabilities}. J. Political Economy.\\
Steele, J. M. (2010). \emph{Stochastic Calculus and Financial Applications}. Springer.
\end{scriptsize}
\end{frame}

\begin{frame}
\frametitle{Black-Scholes. Pricing formula}
\begin{scriptsize}
\begin{itemize}
\item With $\tau = T - t$, log-moneyness $x = \ln(S/K)$ and $u = e^{\alpha x + \beta \tau} v$
for suitable $\alpha, \beta$, the BS PDE reduces to the \textbf{heat equation}
\begin{equation*}
    v_\tau = \tfrac{1}{2} \sigma^2 v_{xx},
    \qquad v(x, 0) = K e^{-\alpha x} (e^x - 1)^+ .
\end{equation*}
\item Convolving with the heat kernel and completing the square gives the \textbf{Black-Scholes formula}
\begin{equation*}
    f(t, S) = S\, \Phi(d_1) - K e^{-r (T - t)} \Phi(d_2),
\end{equation*}
\begin{equation*}
    d_{1,2} = \frac{\ln(S/K) + \left( r \pm \frac{1}{2} \sigma^2 \right) (T - t)}{\sigma \sqrt{T - t}} .
\end{equation*}
\item \textbf{Implied volatility}: the unique $\ivol(K,T) > 0$ such that
$C_{\mathrm{BS}}(\ivol) = C^{\mathrm{mkt}}$.
Uniqueness follows from the vega
\begin{equation*}
    \mathcal{V} = \frac{\partial C_{\mathrm{BS}}}{\partial \sigma} = S \sqrt{T - t}\, \phi(d_1) > 0,
\end{equation*}
and it is computed with Newton's method
$\sigma^{(n+1)} = \sigma^{(n)} - \bigl(C_{\mathrm{BS}}(\sigma^{(n)}) - C^{\mathrm{mkt}}\bigr) / \mathcal{V}(\sigma^{(n)})$.
\end{itemize}
\end{scriptsize}
\end{frame}

\begin{frame}
\frametitle{Black-Scholes. The implied volatility surface}
\begin{scriptsize}
If Black-Scholes described the market, $\ivol(K,T)$ would be flat.
SPY quotes on 30 April 2026 show a \textbf{negative skew} in strike and a \textbf{term structure} in maturity.
\end{scriptsize}
\begin{figure}[htb!]
    \centering
    \begin{subfigure}[b]{0.55\linewidth}\centering
    \includegraphics[width=\linewidth]{iv_composite}
    \caption{\scriptsize Implied volatility surface.}
    \end{subfigure}%
    \begin{subfigure}[b]{0.45\linewidth}\centering
    \includegraphics[width=\linewidth]{iv_smile}
    \caption{\scriptsize Smile at $T \approx 0.21$.}
    \end{subfigure}
\end{figure}
\begin{scriptsize}
Constant $\sigma$ is inconsistent with the market, which motivates the local volatility model.
\end{scriptsize}
\end{frame}
